\section{Related Work}
\label{sec:related}

\paragraph{Indirect Prompt Injection.}
\citet{greshake2023youve} first demonstrated that adversarial instructions embedded in external data sources can hijack LLM agents.
Subsequent work has systematized these attacks: \citet{zhan2024injecagent} introduce InjecAgent with 1{,}054 attack cases across 17 tools, while \citet{debenedetti2024agentdojo} propose AgentDojo, a dynamic environment for evaluating attacks and defenses.
\citet{yi2023benchmarking} and \citet{liu2024formalizing} formalize the threat model and provide attack taxonomies.
More recent work has broadened the attack surface: \citet{shi2025optimization} study optimization-based injection crafting, and \citet{liao2025eia} examine environment-driven injection in agentic settings.
\citet{yuan2024rjudge} and \citet{debenedetti2024agentdojo} highlight the difficulty of automated safety evaluation in dynamic environments.

\paragraph{Existing Benchmarks and Their Limitations.}
Table~\ref{tab:benchmark_comparison} compares \asb{} with existing evaluation frameworks.
InjecAgent provides diverse attack scenarios but evaluates only 1 defense (a basic input classifier).
AgentDojo includes a dynamic execution environment but implements only 2 defenses (spotlighting and output filtering).
ToolEmu \citep{wu2024toolemu} focuses on risk emulation rather than defense evaluation.
None of these benchmarks provides a systematic comparison of diverse defense strategies under controlled conditions.
\asb{} fills this gap by offering: (1)~11 defenses vs.\ 0--2, enabling cross-category comparison; (2)~both ablation and composition experiments to understand why defenses work; (3)~statistical rigor (McNemar's test, Wilson CIs) across 4 models $\times$ 3 runs; and (4)~adaptive attack evaluation.

\begin{table}[t]
\centering
\small
\caption{Comparison of IPI evaluation benchmarks.}
\label{tab:benchmark_comparison}
\resizebox{\columnwidth}{!}{%
\begin{tabular}{lcccc}
\toprule
Feature & InjecAgent & AgentDojo & ToolEmu & \textbf{ASB} \\
\midrule
Attack cases & 1054 & 97 & 36 & \textbf{352} \\
Benign cases & --- & 97 & --- & \textbf{213} \\
Defenses & 1 & 2 & 0 & \textbf{11} \\
Models tested & 1 & 4 & 3 & \textbf{4} \\
Attack types & 4 & 3 & --- & \textbf{6} \\
Composition & \xmark & \xmark & \xmark & \checkmark \\
Ablation & \xmark & \xmark & \xmark & \checkmark \\
Adaptive atk. & \xmark & \checkmark & \xmark & \checkmark \\
Open-source & \checkmark & \checkmark & \checkmark & \checkmark \\
\bottomrule
\end{tabular}%
}
\end{table}

\paragraph{Defense Strategies.}
Proposed defenses include Spotlighting with data marking \citep{hines2024defending}, sandwich prompting \citep{yi2023benchmarking}, re-execution-based verification \citep{xiang2024melon}, policy-based tool constraints \citep{wu2024toolemu}, dual-LLM screening \citep{chen2024struq}, and instruction hierarchy fine-tuning \citep{wallace2024instruction}.
However, each has been evaluated independently on different benchmarks; no prior work provides a controlled comparison under identical conditions---the central contribution of \asb{}.

\paragraph{Contextual Integrity.}
Nissenbaum's theory of contextual integrity \citep{nissenbaum2004privacy,nissenbaum2010privacy} provides a framework for reasoning about appropriate information flows through context-dependent norms.
We apply contextual integrity to LLM agent security in our D10 (\civ{}) defense, which treats the user's goal as defining expected information-flow norms and flags tool calls that violate those norms (Appendix~\ref{sec:civ}).
